\title{DeepSeekMath: Pushing the Limits of Mathematical Reasoning in Open Language Models}

\begin{abstract}
Language models frequently struggle with mathematical reasoning tasks, primarily because of the intricate and highly structured requirements inherent in mathematics. In this work, we present DeepSeekMath~7B, developed by continuing the pre-training of DeepSeek-Coder-Base-v1.5 7B using a corpus of 120B math-specific tokens collected from Common Crawl, which also includes natural language and code samples. DeepSeekMath~7B attains a notable 51.7\% on the competition-grade MATH benchmark, accomplishing this without the use of auxiliary toolkits or ensemble voting approaches, thus drawing near the effectiveness of Gemini-Ultra and GPT-4. Additionally, applying self-consistency over 64 outputs from DeepSeekMath~7B further improves the MATH benchmark result to 60.9\%. The model's advanced mathematical reasoning can be credited to two primary innovations: (1) the construction of a refined data selection pipeline leveraging large-scale, open-access web resources; and (2) the implementation of Group Relative Policy Optimization (GRPO), an adaptation of Proximal Policy Optimization (PPO), which not only advances reasoning skills but also brings greater efficiency to PPO's memory consumption.
\end{abstract}

\section{Introduction}

Large language models (LLMs) have transformed the landscape of mathematical reasoning within artificial intelligence, driving notable progress across benchmarks in quantitative reasoning \citep{MATH} as well as geometry reasoning \citep{trinh2024solving}. These models have also demonstrated value in supporting humans with challenging mathematical tasks \citep{tao}. Despite these strides, the most advanced systems, including GPT-4 \citep{gpt4} and Gemini-Ultra \citep{gemini}, remain inaccessible to the public, while open-source alternatives that are available still lag markedly in terms of effectiveness.

In this work, we present DeepSeekMath, a specialized language model tailored to mathematical content that markedly surpasses other open-source counterparts in mathematical reasoning and nears GPT-4’s benchmark performance on academic tasks. This advancement is made possible through the construction of the DeepSeekMath Corpus, an extensive and high-quality pre-training dataset comprising 120B math tokens. The dataset is sourced from Common Crawl (CC) and filtered using a fastText-based classifier \citep{joulin2016fasttext}. For the classifier’s initial training, we designate data from OpenWebMath \citep{openwebmath} as positive samples, paired with a broad range of alternative web pages representing negatives. With this setup, the classifier identifies further positive instances within CC; these are then meticulously curated through human annotation to increase data reliability. After this refinement, the classifier is retrained on the improved dataset to boost its accuracy. Performance assessments reveal the superior quality of the large-scale dataset, with our base model DeepSeekMath-Base 7B achieving 64.2\% accuracy on GSM8K \citep{gsm8k} and 36.2\% on the advanced MATH dataset \citep{MATH}, outperforming models such as Minerva 540B \citep{minerva}. Additionally, due to the multilingual nature of the DeepSeekMath Corpus, we observe notable improvements on Chinese mathematical evaluations \citep{wei2023cmath,agieval}. We hope our strategies for mathematical data handling provide a foundation for further research, with substantial opportunities remaining for future enhancements.

DeepSeekMath-Base builds upon DeepSeek-Coder-Base-v1.5 7B \citep{deepseek-coder}, given our observation that initializing with a model trained on code yields superior outcomes versus beginning with a standard LLM. In addition, we find that training with mathematical tasks also leads to improvements on the MMLU \citep{mmlu} and BBH benchmarks \citep{bbh}, suggesting that this approach not only strengthens the model’s mathematical proficiency but also enhances its general reasoning performance.

Following the pre-training stage, we perform mathematical instruction tuning on DeepSeekMath-Base utilizing datasets focused on chain-of-thought \citep{cot}, program-of-thought \citep{pot,pal}, and tool-integrated reasoning \citep{tora}. The trained model, DeepSeekMath-Instruct 7B, outperforms every other 7B model and demonstrates performance on par with instruction-tuned open-source models at the 70B scale.

Additionally, we present Group Relative Policy Optimization (GRPO), a novel reinforcement learning (RL) algorithm inspired by Proximal Policy Optimization (PPO) \citep{schulman2017proximal}. Unlike traditional approaches, GRPO eliminates the use of a critic model, opting instead to estimate the baseline directly from group scores. This choice leads to a considerable reduction in training requirements. Leveraging only a selected portion of English instruction tuning data, GRPO substantially surpasses the already strong DeepSeekMath-Instruct, achieving notable gains in both in-domain tasks (GSM8K: 82.9\% $\rightarrow$ 88.2\%, MATH: 46.8\% $\rightarrow$ 51.7\%) and out-of-domain mathematical benchmarks (for instance, CMATH: 84.6\% $\rightarrow$ 88.8\%) during RL training. We further propose a unified framework to interpret various approaches, such as Rejection Sampling Fine-Tuning (RFT) \citep{yuan2023scaling}, Direct Preference Optimization (DPO) \citep{dpo}, PPO, and GRPO. Within this unified framework, it becomes apparent that these methods are best viewed as forms of direct or simplified RL strategies. We perform thorough experimental studies—comparing online versus offline training, outcome versus process-based supervision, single-turn compared to iterative RL, and more—to thoroughly analyze the fundamental aspects of this paradigm. Finally, we discuss the mechanisms behind the improvements brought by our RL approach to instruction-tuned models and outline future research avenues for further refining RL methods within this conceptual framework.

\subsection{Contributions}
We present scalable approaches to mathematical pre-training and further investigate and analyze reinforcement learning methodologies.

\textbf{Math Pre-Training at Scale}

\begin{itemize}[topsep=0pt]
    \item Our study demonstrates that the widely available Common Crawl resource holds substantial mathematically relevant material. Through the application of a carefully crafted data curation process, we assembled the DeepSeekMath Corpus—an extensive dataset containing 120B tokens sourced from web pages selectively screened for mathematical content. This makes it nearly seven times larger than the collection of math web pages used by Minerva \citep{minerva}, and roughly nine times greater in scale compared to the latest OpenWebMath dataset \citep{openwebmath}.

\item DeepSeekMath-Base 7B, our foundational pre-trained model, demonstrates performance on par with Minerva 540B \citep{minerva}. This suggests that parameter count alone does not solely determine mathematical reasoning strength. High-quality pre-training data allows a more compact model to perform at a high level.

\item We present results from our experiments involving math training. Pre-training on code enhances the model's performance on mathematical tasks, irrespective of whether auxiliary tools are used. These results contribute to addressing the longstanding inquiry: \textit{does code training enhance reasoning capabilities?} Our evidence supports the notion that, for mathematical reasoning in particular, it does.

\item While it is standard practice to utilize arXiv papers for training, particularly within mathematics-focused studies, this approach yields no significant gains across any of the mathematical benchmarks evaluated in this work.
\end{itemize}

\textbf{Exploration and Analysis of Reinforcement Learning}

\begin{itemize}[topsep=0pt]
    \item We propose Group Relative Policy Optimization (GRPO), a reinforcement learning algorithm characterized by both efficiency and effectiveness. Unlike Proximal Policy Optimization (PPO), GRPO eliminates the need for a critic model, instead deriving the baseline from group-level scores, thereby greatly decreasing the computational resources required for training.
    \item We show that GRPO yields notable improvements in the performance of our instruction-tuned model DeepSeekMath-Instruct, utilizing only instruction-tuning datasets. In addition, we find that GRPO contributes to better generalization in out-of-domain settings throughout reinforcement learning.
    \item We establish a cohesive framework for interpreting approaches such as RFT, DPO, PPO, and GRPO. Our study further involves comprehensive experimental comparisons—including analyses of online versus offline training, outcome versus process supervision, and single-turn versus iterative reinforcement learning—to systematically dissect the core components of this unified framework.
    \item Leveraging our comprehensive paradigm, we investigate the underlying factors contributing to the success of reinforcement learning and delineate several prospective avenues for advancing the effectiveness of reinforcement learning applied to LLMs.
\end{itemize}

\subsection{Summary of Evaluations and Metrics}

\begin{itemize}[topsep=0pt]
    \item \textbf{English and Chinese Mathematical Reasoning}: Our models are rigorously evaluated using both English and Chinese datasets, spanning mathematical questions from elementary to collegiate difficulty. The English evaluation suite consists of GSM8K \citep{gsm8k}, MATH \citep{MATH}, SAT \citep{llemma}, OCW Courses \citep{minerva}, and MMLU-STEM \citep{mmlu}. For Chinese, we utilize MGSM-zh \citep{mgsm}, CMATH \citep{wei2023cmath}, Gaokao-MathCloze \citep{agieval}, and Gaokao-MathQA \citep{agieval}. Our assessment covers the capacity of models to independently generate comprehensive text-based solutions, as well as their proficiency in employing Python for problem-solving.

On English-language evaluations, DeepSeekMath-Base performs on par with the proprietary Minerva 540B model \citep{minerva}, and outperforms all publicly available base models—including Mistral 7B \citep{mistral} and Llemma-34B \citep{llemma}—by a large margin, irrespective of their exposure to mathematics-specific pre-training. Importantly, DeepSeekMath-Base achieves even greater success on Chinese tasks. This improvement likely stems from our decision to expand beyond the English-only mathematical pre-training datasets utilized in previous research \citep{minerva,llemma}, integrating substantial quantities of high-quality data in other languages. By further applying mathematical instruction tuning and reinforcement learning, the upgraded DeepSeekMath-Instruct and DeepSeekMath-RL models achieve robust results, surpassing the 50\% accuracy threshold on the challenging, competition-oriented MATH dataset—an unprecedented accomplishment for open-source systems.

\item \textbf{Formal Mathematics}: DeepSeekMath-Base is assessed through the informal-to-formal theorem proving benchmark detailed in \citep{dsp_proof}, leveraging the miniF2F dataset \citep{minif2f} with Isabelle \citep{isabelle} serving as the designated proof assistant. The results highlight DeepSeekMath-Base’s robust ability to perform few-shot autoformalization.
\item \textbf{Natural Language Understanding, Reasoning, and Code}: In order to thoroughly characterize the model’s capabilities in general understanding, reasoning, and programming, we subject DeepSeekMath-Base to the Massive Multitask Language Understanding (MMLU) benchmark \citep{mmlu}, which features 57 multiple-choice tasks spanning a wide array of disciplines; BIG-Bench Hard (BBH) \citep{bbh}, a suite of 23 complex tasks emphasizing multi-step reasoning; and two standard code evaluation benchmarks, HumanEval \citep{codex} and MBPP \citep{mbpp}. Mathematical pre-training yields notable improvements in both reasoning and language comprehension tasks.

\section{Math Pre-Training}

\subsection{Data Collection and Decontamination} This section details the methodology used to assemble the DeepSeekMath Corpus utilizing data from Common Crawl. Figure \ref{fig:data_collect} illustrates a cyclical pipeline designed to efficiently accumulate a substantial mathematical text corpus from Common Crawl, beginning with a foundational seed corpus consisting of a limited but high-quality set of mathematics-focused datasets. It should be emphasized that this method is readily transferable to additional fields, including programming.

Initially, we select OpenWebMath \citep{openwebmath}, a curated dataset of mathematically rich web documents, as our primary seed corpus. Using this data, we train a fastText model \citep{joulin2016fasttext} with the objective of retrieving additional web pages with characteristics similar to OpenWebMath. To build our training set, we randomly draw 500,000 instances from the seed corpus for positive samples, while the negative samples consist of 500,000 web pages sourced from Common Crawl. The model training utilizes an open-source toolkit\footnote{\url{https://fasttext.cc}}, with the following hyperparameters: a vector size of 256, a learning rate of 0.1, a maximum n-gram length set to 3, a minimum word count of 3, and a total of 3 epochs. In order to condense the Common Crawl dataset, we implement deduplication at the URL level along with near-deduplication methods, reducing the dataset to 40 billion HTML web pages. We then leverage the trained fastText model to retrieve mathematical web content from this deduplicated set. Next, pages are ranked based on their fastText-generated scores, and only the highest-scoring (and thus presumed highest quality) pages are retained, thereby filtering out less relevant material. We assess how much data to keep by conducting preliminary training runs using the highest-ranking subsets containing 40B, 80B, 120B, and 160B tokens. For this initial cycle, we opt to retain the subset consisting of the top 40B tokens.

Following the initial round of data collection, a significant number of mathematical web pages are left uncollected, primarily due to insufficient diversity among the positive examples used to train the fastText model. To address this issue, we seek out and incorporate additional mathematical web sources to broaden the seed corpus and enhance the model’s effectiveness. Our process begins by partitioning the complete Common Crawl into separate domains, each defined by a unique base URL. For each domain, we compute the fraction of web pages gathered during the first iteration. Domains in which more than 10\% of web pages have been acquired are flagged as math-related (e.g., \url{mathoverflow.net}). Next, we conduct a manual annotation of URLs within these domains that contain mathematical content (such as \url{mathoverflow.net/questions}). Uncollected web pages corresponding to these URLs are subsequently included in the seed corpus. By expanding the set of positive examples, we are able to retrain the fastText model to recognize and retrieve a greater proportion of mathematical data in the following iterations. After completing four rounds of data collection, we have amassed 35.5M mathematical web pages, encompassing 120B tokens. By the fourth iteration, we observe that almost 98\% of the collected data was already obtained during the third iteration, prompting us to discontinue further data acquisition.

In order to prevent benchmark contamination, we adopt the filtering strategy used in \cite{deepseek-coder}, excluding any web content that features questions or answers from English mathematical benchmarks, including GSM8K \citep{gsm8k}, MATH \citep{MATH}, and Chinese benchmarks such as CMATH \citep{wei2023cmath} and AGIEval \citep{agieval}. Our approach involves discarding any text passage from our math training data if it contains a 10-gram sequence that perfectly matches a segment from any of the evaluation benchmarks. For benchmark materials shorter than 10 grams but comprising at least 3 grams, we apply precise string matching to identify and exclude contaminated pages.

\subsection{Validating the Quality of the DeepSeekMath~Corpus}

We run pre-training experiments to investigate how the DeepSeekMath~Corpus is compared with the recently released math-training corpora:

\begin{itemize}[topsep=0pt]
    \item \textbf{MathPile} \citep{mathpile}: a collection combining multiple data sources, amounting to 8.9B tokens, drawn from textbooks, Wikipedia, ProofWiki, CommonCrawl, StackExchange, and arXiv—with arXiv constituting more than 85\% of the dataset;
    \item \textbf{OpenWebMath} \citep{openwebmath}: a 13.6B token dataset built from CommonCrawl by filtering for documents featuring mathematical content;
    \item \textbf{Proof-Pile-2} \citep{llemma}: this corpus merges OpenWebMath, AlgebraicStack (which consists of 10.3B tokens of mathematics-focused code), and arXiv papers (comprising 28.0B tokens). For experiments involving Proof-Pile-2, we employ the arXiv:Web:Code ratio of 2:4:1 as suggested in \cite{llemma}.
\end{itemize}

\subsubsection{Training Setting}
\label{sec:quality-policy}
We conduct mathematics-specific training on a foundational language model comprising 1.3B parameters, structured identically to the DeepSeek LLM architecture \citep{deepseek-llm}, and referred to as DeepSeek-LLM 1.3B. The model undergoes individual training on each distinct mathematical dataset, each for a total of 150B tokens. Our experiments utilize the streamlined HAI-LLM framework \citep{haillm}, chosen for its efficiency and lightweight design. Adhering to methodologies used in DeepSeek LLM development, we employ the AdamW optimizer \citep{adamW}, configured with $\beta_1=0.9$, $\beta_2=0.95$, and a $\mathrm{weight\_decay}$ of 0.1. The learning rate follows a multi-phase schedule: it peaks after 2,000 steps of warming up, is reduced to 31.6\% of its maximum value at the 80\% completion point of training, and is further decayed to 10.0\% by 90\% of training. The peak learning rate is set at 5.3e-4. Each training run uses a batch size of 4M tokens and a context window of 4K tokens.

\subsubsection{Evaluation Results}

\textbf{The DeepSeekMath~Corpus is of high quality, covers multilingual mathematical content, and is the largest in size.}

\begin{itemize}[topsep=0pt]
    \item \textbf{High-quality}:
    To assess the quality of our dataset, we measured downstream effectiveness across eight mathematical benchmarks using few-shot chain-of-thought prompting \cite{cot}. As evidenced in Table \ref{tab:corpora_comparison}, models trained on the DeepSeekMath~Corpus outperformed their counterparts. Additionally, Figure \ref{fig:corpora_comparisons} highlights that, at the 50B token mark (equivalent to one epoch of Proof-Pile-2), models using the DeepSeekMath~Corpus exceeded the performance of those trained with Proof-Pile-2, suggesting a consistently higher average data quality in the DeepSeekMath~Corpus. 
    \item \textbf{Multilingual}:
    The DeepSeekMath~Corpus contains mathematical texts in a variety of languages, with English and Chinese being especially prominent. As seen in Table \ref{tab:corpora_comparison}, using the DeepSeekMath~Corpus in training significantly boosts mathematical reasoning abilities in both English and Chinese. Conversely, other existing corpora, which mainly focus on English, either show limited gains or detract from Chinese reasoning performance.
    \item \textbf{Large-scale}:
    Compared to current mathematical datasets, the DeepSeekMath~Corpus is markedly larger in scale. Figure \ref{fig:corpora_comparisons} shows that DeepSeek-LLM 1.3B, when trained on this corpus, achieves sharper and more sustained improvements. On the other hand, baseline corpora, due to their smaller size and repeated reuse during training, cause model performance to plateau more rapidly.
\end{itemize}

\subsection{Training and Evaluating DeepSeekMath-Base 7B}

This section details the development of DeepSeekMath-Base 7B, a foundational model designed with advanced mathematical reasoning capabilities. Building on DeepSeek-Coder-Base-v1.5 7B \citep{deepseek-coder}, we conducted further training for a total of 500B tokens. The data utilized in this process was distributed as follows: 56\% originated from the DeepSeekMath Corpus, 4\% from AlgebraicStack, 10\% from arXiv, 20\% from Github code repositories, while the final 10\% consisted of natural language sources from Common Crawl in both English and Chinese. Most training configurations follow the protocol in Section \ref{sec:quality-policy}, with two exceptions: the learning rate's upper limit is adjusted to 4.2e-4, and each batch processes up to 10M tokens.

Our evaluation thoroughly examines the mathematical proficiency of DeepSeekMath-Base 7B, emphasizing its capacity to independently generate complete mathematical solutions, utilize external tools for solving mathematical tasks, and perform formal theorem proving. In addition to mathematics, we extend our analysis to encompass a broader overview of the base model's general competencies, covering natural language understanding, reasoning abilities, and programming proficiency.

\paragraph{Mathematical Problem Solving with Step-by-Step Reasoning}

We assess the ability of DeepSeekMath-Base to solve mathematical problems using few-shot chain-of-thought prompting \citep{cot}, conducting experiments on eight different English and Chinese benchmarks. These datasets include tasks that require quantitative reasoning, such as GSM8K \citep{gsm8k}, MATH \citep{MATH}, and CMATH \citep{wei2023cmath}, as well as multiple-choice assessments like MMLU-STEM \citep{mmlu} and Gaokao-MathQA \citep{agieval}. The problems span a broad spectrum of mathematical areas, ranging from elementary school to collegiate levels of difficulty.

According to Table \ref{tab:base_math_cot}, DeepSeekMath-Base 7B consistently achieves the highest scores across all eight benchmarks when compared to open-source base models, which include popular models such as Mistral 7B \citep{mistral} and Llemma 34B \citep{llemma} (the latter extensively trained on Proof-Pile-2 \citep{llemma}). Remarkably, on the challenging, competition-level MATH dataset, DeepSeekMath-Base outperforms all existing open-source base models by a margin exceeding 10\% in absolute terms. Furthermore, it surpasses Minerva 540B \citep{minerva}, a proprietary model based on PaLM \citep{palm}, which is 77 times larger and fine-tuned with mathematical corpora.

\paragraph{Mathematical Problem Solving with Tool Use} We assess the ability of models to perform program-assisted mathematical reasoning on the GSM8K and MATH datasets by employing few-shot program-of-thought prompting \citep{pot,pal}. In this setup, the models are instructed to generate a Python script as a solution, with access to tools such as the \textit{math} and \textit{sympy} libraries to handle complex calculations. The final response is derived by executing the produced code. According to the results presented in Table \ref{tab:base_math_pot_proof}, DeepSeekMath-Base 7B achieves superior performance compared to the previous best model, Llemma 34B.

\paragraph{Formal Mathematics}

The automation of formal proof has become increasingly important for verifying the correctness and robustness of mathematical proofs while also improving productivity. In this work, we assess the capabilities of DeepSeekMath-Base 7B on the informal-to-formal proving task described in \citep{dsp_proof}, where the objective is to produce a formal proof given an informal statement, its formal version, and an informal proof. Our evaluation utilizes the miniF2F benchmark \citep{minif2f}, which targets formalization challenges at the Olympiad level, and involves prompting the model with a few examples to generate formal proofs in Isabelle for each test problem. Consistent with the methodology of \cite{dsp_proof}, our approach relies on the model to output proof sketches, while subsequently employing the general-purpose automated prover Sledgehammer \citep{sledgehammer} to complete any incomplete proof components. Table \ref{tab:base_math_pot_proof} reports that DeepSeekMath-Base 7B attains notable results in the autoformalization of proofs.

\paragraph{Natural Language Understanding, Reasoning, and Code}

We assess the abilities of various models in natural language understanding using MMLU \citep{mmlu}, in reasoning using BBH \citep{bbh}, and in coding tasks with HumanEval \citep{codex} and MBPP \citep{mbpp}. Table \ref{tab:understanding_reasoning_code} indicates that DeepSeekMath-Base 7B achieves notable improvements compared to its predecessor, DeepSeek-Coder-Base-v1.5 \citep{deepseek-coder}, on both the MMLU and BBH benchmarks. This suggests that specialized math-oriented training benefits both language comprehension and reasoning skills. Moreover, through the ongoing incorporation of code tokens in its training process, DeepSeekMath-Base 7B is able to preserve the coding performance demonstrated by DeepSeek-Coder-Base-v1.5 on HumanEval and MBPP. In summary, DeepSeekMath-Base 7B outperforms the general-purpose model Mistral 7B \citep{mistral} across all three benchmarks for reasoning and coding.

\section{Supervised Fine-Tuning}

\subsection{SFT Data Curation}

We develop a comprehensive instruction-tuning dataset for mathematics that encompasses problems in both English and Chinese, drawn from a range of mathematical disciplines and exhibiting diverse levels of difficulty. Each problem is matched with a corresponding solution provided in the chain-of-thought (CoT) \citep{cot}, program-of-thought (PoT) \citep{pot,pal}, and tool-integrated reasoning formats \citep{tora}. Altogether, the dataset contains 776,000 training instances.

\begin{itemize}[topsep=0pt]
    \item \textbf{English mathematical datasets}: We provide tool-assisted annotations for GSM8K and MATH, and include a selected portion of MathInstruct \citep{MathInstruct} as well as the training split from Lila-OOD \citep{lila}, where Chain-of-Thought (CoT) or Program-of-Thought (PoT) approaches are utilized in solutions. This English dataset encompasses a wide spectrum of mathematical disciplines, such as algebra, probability, number theory, calculus, and geometry.
    \item \textbf{Chinese mathematical datasets}: We assemble a set of Chinese K-12 mathematics questions distributed across 76 sub-categories—for instance, linear equations—and furnish solutions leveraging both CoT strategies and tool-based reasoning annotations.
\end{itemize}

\subsection{Training and Evaluating DeepSeekMath-Instruct 7B}

In this section, we present DeepSeekMath-Instruct 7B, a model that is fine-tuned for mathematical instruction using DeepSeekMath-Base as its foundation. Training data are concatenated at random until the maximum context window of 4K tokens is filled. The training process comprises 500 steps, utilizing a batch size of 256 and maintaining a steady learning rate of 5e-5 throughout.

We evaluate models' mathematical performance both without and with tool use, on 4 quantitative reasoning benchmarks in English and Chinese. We benchmark our model against the leading models of the time:

\begin{itemize}[topsep=0pt]
    \item \textbf{Closed-source models} comprise the following: (1) the GPT series, notably GPT-4 \citep{gpt4} and the GPT-4 Code Interpreter~\footnote{\url{https://openai.com/blog/chatgpt-plugins\#\#code-interpreter}}, which currently stand out as the most advanced; (2) Gemini Ultra and Pro \citep{gemini}; (3) Inflection-2 \citep{inflection-2}; (4) Grok-1~\footnote{\url{https://x.ai/model-card}}; and a number of recent models introduced by Chinese firms, such as (5) Baichuan-3~\footnote{\url{https://www.baichuan-ai.com}} and (6) the latest GLM-4~\footnote{\url{https://open.bigmodel.cn/dev/api\#glm-4}}.

    from the GLM family \citep{glm}.

These models are designed for broad usage and typically have undergone multiple alignment steps.  
\item \textbf{Open-source models} consist of: general-purpose models such as (1) DeepSeek-LLM-Chat 67B \citep{deepseek-llm}, (2) Qwen 72B \citep{qwen}, (3) SeaLLM-v2 7B \citep{seallm}, and (4) ChatGLM3 6B \citep{chatglm3}. In addition, several models include mathematics-specific enhancements: (5) InternLM2-Math 20B~\footnote{\url{https://github.com/InternLM/InternLM-Math}}, which extends InternLM2 by training on mathematical tasks and then performing instruction tuning; (6) Math-Shepherd-Mistral 7B, which leverages PPO training \citep{schulman2017proximal} with a process-supervised reward model on top of Mistral 7B \citep{mistral}; (7) the WizardMath series \citep{wizardmath}, which strengthens mathematical reasoning in Mistral 7B and Llama-2 70B \citep{llama2} via evolve-instruct (an instruction tuning approach using instructions generated by AI) and PPO, primarily with training data from GSM8K and MATH; (8) MetaMath 70B \citep{metamath}, formed by fine-tuning Llama-2 70B on an expanded version of the GSM8K and MATH datasets; (9) ToRA 34B \cite{tora}, in which CodeLlama 34B is further tuned for mathematical reasoning augmented with tool use; (10) MAmmoTH 70B \citep{MathInstruct}, created by instruction-tuning Llama-2 70B on MathInstruct.

Table \ref{tab:sft_rl_math} illustrates that in scenarios where tool use is prohibited, DeepSeekMath-Instruct 7B achieves impressive step-by-step reasoning capabilities. Specifically, on the competition-level MATH dataset, our model outperforms all open-source contenders and most commercial models (including Inflection-2 and Gemini Pro) by a margin of no less than 9\% absolute. This performance edge extends even against models that are much larger in scale (such as Qwen 72B) or have undergone specialized math-oriented reinforcement learning (for example, WizardMath-v1.1 7B). Although DeepSeekMath-Instruct reaches parity with Chinese proprietary systems like GLM-4 and Baichuan-3 on the MATH benchmark, it does not yet reach the level of GPT-4 or Gemini Ultra.

When evaluated in scenarios that permit the combination of natural language reasoning with programmatic tool usage to tackle problems, DeepSeekMath-Instruct 7B achieves nearly 60\% accuracy on the MATH benchmark, outperforming every other available open-source model. Furthermore, across additional benchmarks, our model demonstrates performance on par with DeepSeek-LLM-Chat 67B—the previous leading model, which boasts ten times the parameter count.

\subsection{Group Relative Policy Optimization} After the Supervised Fine-Tuning (SFT) phase, reinforcement learning (RL) has demonstrated significant success in enhancing the mathematical reasoning skills of LLMs \citep{wang2023math,wizardmath}. Here, we present Group Relative Policy Optimization (GRPO), our novel RL approach designed for both efficiency and effectiveness.

\subsubsection{From PPO to GRPO} Proximal Policy Optimization (PPO) \citep{schulman2017proximal} is an actor-critic RL algorithm that is widely used in the RL fine-tuning stage of LLMs \citep{ouyang2022training}. In particular, it optimizes LLMs by maximizing the following surrogate objective:

\begin{equation}
\footnotesize
    \mathcal{J}_{PPO}(\theta) = \mathbb{E}{[q \sim P(Q),\ o \sim \pi_{\theta_{old}}(O|q)]} \frac{1}{|o|} \sum_{t=1}^{|o|} \min \left[ \frac{\pi_\theta(o_{t} | q, o_{<t})}{\pi_{\theta_{old}}(o_{t} | q, o_{<t})} A_{t},\ \text{clip}\left(\frac{\pi_\theta(o_{t} | q, o_{<t})}{\pi_{\theta_{old}}(o_{t} | q, o_{<t})},\ 1 - \varepsilon,\ 1 + \varepsilon \right) A_{t} \right] ,
\end{equation}

This equation represents the PPO objective function $\mathcal{J}_{PPO}(\theta)$, where the expectation is computed over samples $q$ drawn from $P(Q)$ and output sequences $o$ sampled from the previous policy $\pi_{\theta_{old}}(O|q)$. The average across each timestep $t$ in the sequence $o$ involves taking the minimum between the importance sampling ratio multiplied by the advantage $A_t$, and the same quantity but with the ratio clipped to the interval $[1 - \varepsilon, 1 + \varepsilon]$. This form helps prevent excessively large updates and thus stabilizes policy optimization.

Here, $\pi_{\theta}$ denotes the current policy model, while $\pi_{\theta_{old}}$ represents the previous policy iteration. The variables $q$ and $o$ correspond to questions drawn from the dataset and outputs generated by the prior policy $\pi_{\theta_{old}}$, respectively. The hyper-parameter $\varepsilon$, which governs the clipping range, is utilized in PPO to enhance the reliability of the training process. The term $A_t$ stands for the advantage estimate, calculated through the Generalized Advantage Estimation (GAE) method \citep{gae}, utilizing future rewards $\{r_{\ge t}\}$ in conjunction with the learned value function $V_{\psi}$. Consequently, PPO requires concurrent training of a value function with the policy network. To reduce excessive reward model optimization, it is customary to include a KL divergence penalty from a reference model into the per-token reward \citep{ouyang2022training}, as described by the following expression:

\begin{equation}
    r_{t} = r_\varphi(q, o_{\le t}) - \beta \log\frac{\pi_{\theta}(o_{t}|q, o_{<t})}{\pi_{ref}(o_{t}|q, o_{<t})},
\label{eq:PPO-reward}
\end{equation}

Here, $r_\varphi$ represents the reward model, $\pi_{ref}$ denotes the reference model—commonly set as the initial SFT model—and $\beta$ indicates the strength of the KL penalty term.

As the value function employed in PPO is typically another model of comparable size as the policy model, it brings a substantial memory and computational burden. Additionally, during RL training, the value function is treated as a baseline in the calculation of the advantage for variance reduction. While in the LLM context, usually only the last token is assigned a reward score by the reward model, which may complicate the training of a value function that is accurate at each token. To address this, as shown in Figure \ref{fig:grpo}, we propose Group Relative Policy Optimization (GRPO), which obviates the need for additional value function approximation as in PPO, and instead uses the average reward of multiple sampled outputs, produced in response to the same question, as the baseline. More specifically, for each question $q$, GRPO samples a group of outputs $\{o_1, o_2, \cdots, o_G\}$  from the old policy  $\pi_{\theta_{old}}$  and then optimizes the policy model by maximizing the following objective:

\begin{equation}
\footnotesize
\begin{split}
    \mathcal{J}_{GRPO}(\theta) &= \mathbb{E}_{q \sim P(Q), \{o_i\}_{i=1}^G \sim \pi_{\theta_{old}}(O|q)} \Bigg[ \\
    & \quad \frac{1}{G} \sum_{i=1}^G \frac{1}{|o_i|} \sum_{t=1}^{|o_i|} \Bigg( \min \Bigg( \frac{\pi_\theta(o_{i,t} \mid q, o_{i,<t})}{\pi_{\theta_{old}}(o_{i,t} \mid q, o_{i,<t})} \hat{A}_{i,t}, \\
    & \quad\quad\quad\quad\quad \text{clip} \left( \frac{\pi_\theta(o_{i,t} \mid q, o_{i,<t})}{\pi_{\theta_{old}}(o_{i,t} \mid q, o_{i,<t})}, 1-\varepsilon, 1+\varepsilon \right) \hat{A}_{i,t} \Bigg) \\
    & \quad\quad\quad\quad - \beta \mathbb{D}_{KL}\left[\pi_{\theta} \Vert \pi_{ref}\right] \Bigg) \Bigg],
\end{split}
\label{eq:GRPO-obj}
\end{equation}

where $\varepsilon$ and $\beta$ are hyper-parameters, and $\hat{A}_{i,t}$  is the advantage calculated based on relative rewards of the outputs inside each group only, which will be detailed in the following subsections. The group relative way that GRPO leverages to calculate the advantages, aligns well with the comparative nature of rewards models,  as reward models are typically trained on datasets of comparisons between outputs on the same question. Also note that, instead of adding KL penalty in the reward, GRPO regularizes by directly adding the KL divergence between the trained policy and the reference policy to the loss, avoiding complicating the calculation of  $\hat{A}_{i,t}$. And different from the KL penalty term used in (\ref{eq:PPO-reward}), we estimate the KL divergence with the following unbiased estimator \citep{kl_approx}:

\begin{equation}
\small
    \mathbb{D}_{KL}\left[\pi_{\theta} || \pi_{ref}\right] = \frac{\pi_{ref}(o_{i,t}|q,o_{i,<t})}{\pi_{\theta}(o_{i,t}|q,o_{i,<t})} - \log\left(\frac{\pi_{ref}(o_{i,t}|q,o_{i,<t})}{\pi_{\theta}(o_{i,t}|q,o_{i,<t})}\right) - 1,
\end{equation}

which is guaranteed to be positive.

\begin{algorithm}[t]
  \small
  \caption{Iterative Group Relative Policy Optimization}
  \textbf{Input} starting policy model $\pi_{\theta_{\text{init}}}$; reward evaluators $r_\varphi$; collection of task prompts $\mathcal{D}$; 
  hyperparameters $\varepsilon$, $\beta$, $\mu$
  \begin{algorithmic}[1]
    \State Initialize policy model $\pi_\theta \gets \pi_{\theta_{\text{init}}}$
    \For{iteration = 1 to $I$}
       \State Assign $\pi_{ref} \gets \pi_{\theta}$
      \For{step = 1 to $M$}
      \State Randomly select mini-batch $\mathcal{D}_b$ from $\mathcal{D}$
      \State Set previous policy $\pi_{\theta_{old}} \gets \pi_{\theta}$ 
      \State For every question $q$ in $\mathcal{D}_b$, draw $G$ samples $\{o_i\}_{i=1}^G \sim \pi_{\theta_{old}} (\cdot \mid q) $
      \State Determine rewards $\{r_i\}_{i=1}^{G}$ for outputs $o_i$ using $r_\varphi$
      \State Calculate group relative advantages $\hat{A}_{i,t}$ for the $t$-th token in each $o_i$
      \For{GRPO iteration = 1 to $\mu$}
        \State Improve policy model $\pi_\theta$ by maximizing the GRPO objective (Equation \ref{eq:GC-GRPO})
      \EndFor
    \EndFor 
    \State Perform continual updates to $r_\varphi$ leveraging a replay buffer. 
    \EndFor 
  \end{algorithmic}
  \textbf{Output} $\pi_\theta$
  \label{alg:iter-grpo}
\end{algorithm}

\subsubsection{Outcome Supervision RL with GRPO} More precisely, given a question $q$, a collection of responses $\{o_1, o_2, \cdots, o_G\}$ is generated by sampling from the previous policy $\pi_{\theta_{old}}$. Each response is evaluated by a reward model, resulting in $G$ reward values $\mathbf{r} = \{r_1, r_2, \cdots, r_G\}$. These reward scores are then standardized by first subtracting the mean of the group and subsequently dividing by the standard deviation. In outcome supervision, the standardized reward for each generated response $o_i$ is used as the terminal reward, which is assigned to all token-wise advantages $\hat{A}_{i, t}$ in the output, i.e., $\hat{A}_{i, t} = \widetilde{r}_i = \frac{r_i- {\rm mean}(\mathbf{r})}{{\rm std}(\mathbf{r})}$. The policy is then improved by maximizing the objective outlined in equation (\ref{eq:GRPO-obj}).

\subsubsection{Process Supervision RL with GRPO}
Relying solely on outcome supervision, which grants a reward only upon completion of the output, may be inadequate and suboptimal when directing the policy in the context of intricate mathematical reasoning. In accordance with \cite{wang2023math}, our investigation also includes process supervision, which allocates rewards after every individual reasoning step. More specifically, for a given question $q$ and a set of $G$ sampled outputs $\{o_1, o_2, \cdots, o_G\}$, a process reward model assigns a score to each reasoning step within the outputs, resulting in a collection of rewards: $\mathbf{R} = \{ \{r_1^{index(1)},\cdots,r_1^{index(K_1)}\}, \cdots,  \{r_G^{index(1)},\cdots,r_G^{index(K_G)}\} \}$. Here, $index(j)$ indicates the position of the terminal token for the $j$-th reasoning step, and $K_i$ denotes how many steps are present in the $i$-th generated output. To ensure consistency and comparability, these reward values are standardized using the overall mean and standard deviation: $\widetilde{r}_i^{index(j)} = \frac{r_i^{index(j)} - {\rm mean(\mathbf{R})}}{{\rm std(\mathbf{R})}}$.

Following this, the process supervision approach defines the advantage at each token as the cumulative sum of the normalized rewards corresponding to all subsequent steps, specifically: $\hat{A}_{i, t} = \sum_{index(j) \ge t} \widetilde{r}_i^{index(j)}$.
Policy optimization is then carried out by maximizing the objective function as expressed in equation (\ref{eq:GRPO-obj}).

\subsubsection{Iterative RL with GRPO} During reinforcement learning, the previously trained reward model might become inadequate for effectively guiding the current policy model. To address this, we investigate an iterative approach to RL using GRPO. As outlined in Algorithm \ref{alg:iter-grpo}, this method involves generating fresh training datasets for the reward model by sampling outputs from the updated policy model. The reward model is then updated iteratively, employing a replay strategy that retains 10\% of past data. Subsequently, the latest policy model is used as the reference model, and training of the policy model proceeds with feedback from the updated reward model.

\subsection{Training and Evaluating DeepSeekMath-RL}

Our reinforcement learning (RL) experiments utilize DeepSeekMath-Instruct 7B as the policy model. The RL training dataset comprises approximately 144K chain-of-thought questions derived from the SFT data, specifically those linked to GSM8K and MATH, while omitting other SFT instances. This selective approach is employed to assess how RL affects benchmarks that lack coverage during the RL stage. Reward model training data are curated in accordance with the method described in \citep{wang2023math}. The initial reward model is based on DeepSeekMath-Base 7B and is trained with a learning rate of $2 \times 10^{-5}$. For the GRPO algorithm, we configure the policy model’s learning rate to $1 \times 10^{-6}$ and apply a KL penalty coefficient of 0.04. For every question, we generate $64$ samples. The maximum sequence length is limited to 1024, and the batch size per training iteration is also 1024. Following each exploration phase, the policy model is updated only once. DeepSeekMath-RL 7B is evaluated on benchmarks as per the protocols used for DeepSeekMath-Instruct 7B. Within DeepSeekMath-RL 7B, tasks utilizing GSM8K and MATH with chain-of-thought problem-solving are categorized as in-domain, while all remaining benchmarks are classified as out-of-domain.

Table \ref{tab:sft_rl_math} presents the results achieved by both open- and closed-source models employing chain-of-thought and tool-augmented reasoning approaches on English and Chinese evaluation datasets. The key observations are as follows: 1) Utilizing chain-of-thought reasoning, DeepSeekMath-RL 7B achieves 88.2\% accuracy on GSM8K and 51.7\% on MATH, outperforming every open-source model within the 7B–70B range and exceeding the capabilities of most closed-source systems. 2) Notably, DeepSeekMath-RL 7B is exclusively trained on instruction tuning data in the chain-of-thought format for GSM8K and MATH, initializing from DeepSeekMath-Instruct 7B. Even with its narrow training data scope, it consistently surpasses DeepSeekMath-Instruct 7B across all metrics, highlighting the advantages provided by reinforcement learning.

\section{Discussion}
This section presents an overview of the results obtained from our pre-training and reinforcement learning experiments.

\subsection{Lessons Learnt in Pre-Training}

We begin by detailing our pre-training process. Except where indicated otherwise, the experiments follow the training configuration described in Section \ref{sec:quality-policy}. Notably, references to the DeepSeekMath Corpus in this discussion pertain to an 89B-token dataset, derived from the second round of data collection.

\subsubsection{Code Training Benefits Mathematical Reasoning}

An often-cited but unproven theory claims that practicing coding can enhance reasoning skills. In this work, we seek to provide partial evidence regarding this claim in the context of mathematics: code training enhances models' mathematical reasoning capabilities, regardless of whether external tools are employed.

To study how code training affects mathematical reasoning, we experimented with the following two-stage training and one-stage training settings:

\textbf{Two-Stage Training}

\begin{itemize}[topsep=0pt]
    \item \textbf{Code Training for 400B Tokens $\rightarrow$ Math Training for 150B Tokens}: The DeepSeek-LLM 1.3B model undergoes pre-training on 400B tokens consisting of code, which is subsequently followed by an additional 150B tokens from mathematical datasets;
    \item \textbf{General Training for 400B Tokens $\rightarrow$ Math Training for 150B Tokens}: For comparison, we conduct a baseline experiment in which the initial 400B tokens comprise general-domain data (sourced from a comprehensive general corpus curated by DeepSeek-AI) rather than code, allowing us to explore how pre-training on code, as opposed to general tokens, influences the model’s ability to reason about mathematical tasks.
\end{itemize}

\textbf{One-Stage Training}

\begin{itemize}[topsep=0pt]
    \item \textbf{Math Pretraining with 150B Tokens}: DeepSeek-LLM 1.3B undergoes pretraining utilizing 150B tokens sourced from mathematical data;
    \item \textbf{Joint Pretraining on 400B Code Tokens and 150B Math Tokens}: We observe that sequential math pretraining after code pretraining diminishes the model's coding ability. Therefore, we assess whether simultaneously pretraining on a combination of code and math tokens can foster mathematical reasoning while helping prevent catastrophic forgetting of coding skills.
\end{itemize}

\paragraph{Results}
The outcomes presented in Table \ref{tab:code-math} and Table \ref{tab:code-math-general-eval} illustrate how various training configurations impact downstream task performance.

Incorporating code training improves program-aided mathematical reasoning, regardless of whether it is applied within a two-stage or a one-stage training framework. As indicated by Table \ref{tab:code-math}, within the two-stage setup, code training by itself notably boosts performance on solving GSM8K and MATH problems using Python. Introducing math training in a subsequent stage further enhances these results. Notably, in the one-stage setting, jointly presenting code tokens and math tokens not only prevents the catastrophic forgetting observed in the two-stage approach, but also produces complementary gains for both coding proficiency (Table \ref{tab:code-math-general-eval}) and program-aided mathematical reasoning (Table \ref{tab:code-math}).

Training on code additionally enhances mathematical reasoning abilities even in the absence of tool assistance. Within the two-stage training framework, the preliminary phase of code training yields notable improvements on its own. This stage also facilitates more effective and accelerated learning during subsequent mathematical training, culminating in superior overall outcomes. Conversely, when code tokens and math tokens are combined in a single-stage training approach, performance in mathematical reasoning without tool use deteriorates. A possible explanation is that DeepSeek-LLM 1.3B’s relatively small model size does not afford it the capacity to comprehensively integrate both code and mathematical data within a single training stage.

\subsubsection{ArXiv Papers Seem Ineffective in Improving Mathematical Reasoning}

ArXiv papers are commonly included as a component of math pre-training data \citep{minerva,gpt-f,llemma,mathpile}. However, detailed analysis regarding their impact on mathematical reasoning has not been extensively conducted. Perhaps counter-intuitively, according to our experiments, arXiv papers seem ineffective in improving mathematical reasoning. We experiment with models of different sizes, including DeepSeek-LLM 1.3B and DeepSeek-Coder-Base-v1.5 7B \citep{deepseek-coder}, using arXiv corpora that underwent varied processing pipelines:

\begin{itemize}[topsep=0pt]
    \item \textbf{MathPile} \citep{mathpile}: a dataset comprising 8.9 billion tokens, curated through a set of cleaning and filtering heuristics, with scientific arXiv papers constituting more than 85\% of its content;
    \item \textbf{ArXiv-RedPajama} \citep{redpajama}: consists of all arXiv LaTeX files with extraneous elements such as preambles, comments, macros, and bibliographies excluded, amounting to 28.0 billion tokens in total.
\end{itemize}

During our experiments, we independently pretrain DeepSeek-LLM 1.3B with 150B tokens and DeepSeek-Coder-Base-v1.5 7B with 40B tokens using each arXiv dataset. Our findings indicate that utilizing arXiv papers does not significantly enhance mathematical reasoning. When models are trained solely on arXiv-related corpora, neither exhibits meaningful gains and may even experience declines across several mathematical evaluation tasks of varying complexity considered in this work. These tasks span quantitative reasoning datasets such as GSM8K and MATH (refer to Table \ref{tab:arxiv-cot}), multiple-choice assessments like MMLU-STEM (Table \ref{tab:arxiv-cot}), as well as domains of formal mathematics including miniF2F (Table \ref{tab:arxiv_atp}).

However, this conclusion has its limitations and should be taken with a grain of salt. We have not yet studied:

\begin{itemize}[topsep=0pt]
    \item How arXiv tokens influence certain mathematics-oriented tasks beyond the scope of this work, for instance, the informalization of theorems, which involves translating formal mathematical statements or proofs into more casual or non-technical language;
    \item The potential interaction effects between arXiv tokens and additional data sources when used together;
    \item If scaling models up would reveal further advantages arising from incorporating arXiv papers.
\end{itemize}

Thus, further exploration is required, which we leave for future studies.

\subsection{Insights of Reinforcement Learning}
\subsubsection{Towards to a Unified Paradigm} Here, we introduce a comprehensive framework to systematically examine various training techniques, including SFT, RFT, DPO, PPO, and GRPO. We also perform empirical studies to investigate key components within this framework. In most cases, the gradient of the objective with respect to the parameter $\theta$ for a given training approach can be expressed as follows:

\begin{equation}
    \nabla_{\theta}\mathcal{J}_{\textcolor{red}{\mathcal{A}}}(\theta) = \mathbb{E}[\underbrace{(q,o) \sim \textcolor{red}{\mathcal{D}}}_{Data \ Source}]\left( \frac{1}{|o|} \sum_{t=1}^{|o|}  \underbrace{GC_{{\mathcal{A}}}(q, o, t, \textcolor{red}{\pi_{{rf}}})}_{Gradient \ Coefficient}  \nabla_{\theta}\log \pi_{\theta}(o_t | q, o_{<t})\right).
\label{eq:objective}
\end{equation}

There exist three key components: 1) \textit{Data Source $\mathcal{D}$}, which determines the training data; 2) \textit{Reward Function $\pi_{{rf}}$}, which is the source of the training reward signal; 3) \textit{Algorithm $\mathcal{A}$}: which processes the training data and the reward signal to the gradient coefficient $GC$ that determines the magnitude of the penalty or reinforcement for the data. We analyze several representative methods based on such a unified paradigm:

\begin{itemize}
    \item  \textbf{Supervised Fine-tuning (SFT)}: In SFT, a pretrained model undergoes additional training using curated SFT data selected by humans.
    \item \textbf{Rejection Sampling Fine-tuning (RFT)}: RFT enhances the SFT model by further fine-tuning it with a subset of responses generated from the SFT model for the SFT questions, retaining only those outputs that are validated as correct.
    \item \textbf{Direct Preference Optimization (DPO)}: DPO extends the fine-tuning of the SFT model with extra sampled responses, leveraging a pairwise DPO loss function to guide refinement.
    \item \textbf{Online Rejection Sampling Fine-tuning (Online RFT)}: Unlike RFT, Online RFT begins with the SFT model as the initial policy and continues to train it using outputs dynamically sampled from the evolving policy itself.
    \item \textbf{PPO/GRPO}: In this approach, the policy model starts from the SFT-trained model and is subsequently optimized through reinforcement using outputs sampled from the current policy during training.
\end{itemize}

Table \ref{tab:data_policy} presents an overview of the constituent elements of these methods. For an expanded explanation of the derivation steps, see Appendix \ref{app:analysis-rl}.

\paragraph{Observation about Data Source} We categorize data sources as either online sampling or offline sampling. In the context of online sampling, training data is generated through real-time interactions of the training policy model during exploration. Conversely, offline sampling refers to gathering training data from outputs produced by the original SFT model. Specifically, RFT and DPO utilize offline sampling approaches, whereas Online RFT and GRPO adopt an online sampling strategy.

As illustrated in Figure \ref{fig:rl-analysis}, our results indicate that Online RFT markedly surpasses RFT across two benchmark tasks. In particular, while the performance of Online RFT aligns closely with that of RFT during the initial phases of training, Online RFT demonstrates a clear and pronounced benefit as training progresses, highlighting the effectiveness of the online training approach. This outcome is expected; at the outset, both the actor and the SFT model are highly similar, and the sampled data reflects only subtle distinctions. As training continues, greater divergences emerge in the data produced by the actor, and the dynamic nature of real-time data sampling becomes increasingly advantageous.

\paragraph{Observation about Gradient Coefficient} The algorithm utilizes the reward signal to calculate the gradient coefficient for updating the model’s parameters. In our experiments, we decompose the reward function into two components: `Rule' and `Model'. The `Rule' component evaluates responses based on the factual accuracy of the answers, while the `Model' component relies on a trained reward model that assigns scores to responses; the training data for this reward model are derived from the judgments made under the `Rule' criterion. As demonstrated in Equations \ref{eq:GC-RFT} and \ref{eq:GC-GRPO}, an important distinction exists between GRPO and Online RFT. Specifically, GRPO modifies its gradient coefficient in accordance with the reward values assigned by the reward model, thereby enabling tailored reinforcement or punishment according to the degree of correctness. In comparison, Online RFT does not provide this flexibility; it treats all correct responses the same—applying the same positive reinforcement—and offers no penalization for incorrect answers.

Figure \ref{fig:rl-analysis} illustrates that GRPO outperforms online RFT, underscoring the effectiveness of modifying both positive and negative gradient coefficients. Moreover, GRPO+PS achieves better outcomes than GRPO+OS, suggesting that incorporating step-aware, fine-grained gradient coefficients is advantageous. Additionally, we investigate iterative RL by executing two iterative cycles in our experiments. As depicted in Figure \ref{fig:iter-rl}, we observe that this iterative approach leads to substantial performance gains, with the most pronounced improvement occurring during the initial iteration.

\subsubsection{Why RL Works?} In this study, we utilize reinforcement learning on a selected portion of the instruction tuning data, resulting in substantial performance improvements over the baseline instruction tuning model. To gain further insight into the underlying mechanisms behind this improvement, we compare the Pass@K and Maj@K accuracies between the Instruct and RL models using two benchmark datasets. As illustrated in Figure \ref{fig:maj-pass}, while reinforcement learning leads to noticeable gains in Maj@K, it does not improve Pass@K. These observations suggest that the benefits of reinforcement learning are due primarily to strengthening the model's output distribution—specifically, \textbf{the increased effectiveness derives from amplifying the likelihood of correct answers within the TopK responses, rather than from advancing the model’s core reasoning ability.} Consistent with this, \citep{wang2023making} previously observed a \textbf{misalignment problem} in reasoning tasks with SFT models and demonstrated that the reasoning abilities of these models could be elevated using a range of preference alignment methods \citep{yuan2023rrhf,song2023preference,wang2023making}.

\subsubsection{How to Achieve More Effective RL?}
\label{sec:effec_rl}
Our results indicate that RL is highly effective for mathematical reasoning problems. Additionally, we establish a coherent framework that accommodates a variety of notable training strategies. In this unified view, every method can be interpreted as an instantiation—either as direct or streamlined forms—of RL methodologies. As highlighted in Equation \ref{eq:objective}, three critical factors are involved: Data Source, Algorithm, and Reward Function. We outline several possible avenues for future work concerning these elements.

\paragraph{Data Source} The foundation of all training methods lies in the data source. Within the framework of RL, this refers specifically to unlabeled questions paired with responses generated from the policy model. For this study, our data is drawn exclusively from questions originating during the instruction tuning phase, and corresponding outputs are generated using simple nucleus sampling. We attribute our RL pipeline's selective improvement on Maj@K metrics to these choices. Moving forward, we aim to extend our RL approach to cover question prompts beyond the original distribution, as well as to integrate \textbf{advanced sampling (decoding) strategies}, such as those inspired by tree-search methodologies \citep{yao2023tree}. Furthermore, the adoption of \textbf{efficient inference techniques} \citep{xia-etal-2023-speculative,leviathan2023fast,kwon2023efficient,xia2024unlocking}, which critically influence how effectively policy models can be explored, will be a key consideration in future work.

\paragraph{Algorithms} Algorithms utilize both data and reward signals to adjust gradient coefficients, thereby updating model parameters. According to Equation \ref{eq:objective}, contemporary approaches fundamentally \textbf{TRUST} the reward function’s signal for modulating the conditional probability assigned to particular tokens—either increasing or decreasing it. Nonetheless, the reliability of the reward signal cannot always be guaranteed, particularly for highly intricate tasks. For instance, the PRM800K dataset \citep{lightman2023let}, which underwent detailed annotation by highly trained human annotators, still has an estimated annotation error rate of about 20\%\footnote{\url{https://github.com/openai/prm800k/issues/12\#issuecomment-1728491852}}. Consequently, we aim to investigate reinforcement learning algorithms that demonstrate resilience in the face of noisy reward signals. Our perspective is that alignment techniques based on \textbf{WEAK-TO-STRONG} \citep{burns2023weak} principles have the potential to fundamentally transform prevailing learning algorithm paradigms.

\paragraph{Reward Function} The reward function provides the core training signal in reinforcement learning. Typically, in RL, this signal is produced by a neural reward model. We identify three major avenues for advancing reward models: 1) \textbf{Improving the reward model's generalization capacity.} It is critical for the reward model to generalize effectively to unseen questions and sophisticated decoding results. Failure to do so may lead RL to simply reinforce the existing distribution of LLMs without driving substantive improvements; 2) \textbf{Capturing and representing the reward model's uncertainty.} Incorporating uncertainty measures may bridge the gap between limited reward models and algorithms designed to transition from weak to strong performance; 3) \textbf{Efficient development of precise process reward models} that deliver granular supervision signals during reasoning tasks \citep{lightman2023let,wang2023math}.

\section{Conclusion, Limitation, and Future Work} In this work, we introduce DeepSeekMath, which surpasses all currently available open-source models on the MATH competition benchmark and achieves results close to those of proprietary systems. DeepSeekMath~is based on DeepSeek-Coder-v1.5 7B and is continually trained on 500B tokens, with 120B math-specific tokens predominantly gathered from Common Crawl. Through a comprehensive ablation analysis, we find that web pages are a rich source of high-quality mathematical material, while arXiv contributes less significantly than anticipated. We propose Group Relative Policy Optimization (GRPO), an adaptation of Proximal Policy Optimization (PPO), which enhances mathematical reasoning with reduced memory overhead. Experimental evidence demonstrates that GRPO yields improvements, even when DeepSeekMath-Instruct 7B already performs strongly on established metrics. Furthermore, we offer a unifying framework for understanding various reinforcement learning methods and highlight several promising research avenues for advancing this area.

While DeepSeekMath~performs remarkably well on quantitative reasoning assessments, its proficiency in areas such as geometry and theorem-proving remains less competitive compared to proprietary systems. For example, our preliminary evaluation revealed that the model struggles with triangle and ellipse problems, which suggests possible biases in the selection of pre-training and fine-tuning data. Additionally, due to limitations in model size, DeepSeekMath~demonstrates inferior few-shot learning capabilities compared to GPT-4; whereas GPT-4 benefits substantially from few-shot prompts, DeepSeekMath~exhibits comparable outcomes in both zero-shot and few-shot scenarios. Moving forward, we aim to refine our data selection methodology to develop a higher-quality pre-training corpus. Furthermore, we intend to investigate more effective reinforcement learning strategies for LLMs, as outlined in Section \ref{sec:effec_rl}.

\end{document}